\documentclass[11pt]{article}
\usepackage[a4paper,margin=1in]{geometry}
\usepackage{fontspec}
\usepackage[boldfont=false]{xeCJK}
\setCJKmainfont{FandolSong-Regular.otf}
\usepackage{amsmath}
\usepackage{amssymb}
\usepackage{array}
\usepackage{booktabs}
\usepackage{graphicx}
\usepackage{adjustbox}
\usepackage{enumitem}
\setlist[itemize]{topsep=2pt,partopsep=0pt,parsep=0pt,itemsep=2pt,leftmargin=1.4em}
\setlist[enumerate]{topsep=2pt,partopsep=0pt,parsep=0pt,itemsep=2pt,leftmargin=1.6em}
\usepackage{parskip}
\usepackage{fvextra}
\fvset{breaklines=true,breakanywhere=true,fontsize=\small}
\setlength{\parindent}{0pt}

\begin{document}

\begin{center}
  {\LARGE\bfseries Work, energy and power}\\[0.35em]
  {\large A Level Physics}
\end{center}
\vspace{0.6em}

\section*{Work, energy and power}

\subsection*{Work done by a force}

\textbf{Work} 功 is done when a force moves its point of contact along the line of the force. The work done by a constant force $F$ that causes a \textbf{displacement} 位移 $s$ is

\[
W = F \cdot s \cdot \cos\theta,
\]

where $\theta$ is the angle between the force and the displacement. Only the \textbf{component} 分量 of the force along the displacement does work.

\par\medskip
\noindent\begin{minipage}{\linewidth}
\centering
\adjustbox{max width=\linewidth, max totalheight=0.42\textheight}{%
  \includegraphics{/Users/joschka/Desktop/Programming/eXercise/output/eXam/handouts/a_level_physics/assets/05-p104-work-force-at-angle.pdf}}\\[0.3em]
\emph{Only the component of the force along the displacement (F cos θ) does work}
\end{minipage}\par\medskip

Unit: $\text{J} = \text{N m}$. Work is a \textbf{scalar} 标量.

Special cases:

\begin{itemize}
\item force in the same direction as the motion ($\theta = 0$): $W = Fs$, positive work, \textbf{energy} 能量 given \textbf{to} the object.

\item force at right angles to the motion ($\theta = 90°$): $W = 0$. The \textbf{normal contact force} 支持力 on a car on a flat road does no work.

\item force opposite to the motion ($\theta = 180°$): $W = -Fs$, negative work, energy taken \textbf{from} the object (for example \textbf{friction} 摩擦力).

\end{itemize}

\par\medskip
\noindent\begin{minipage}{\linewidth}
\centering
\adjustbox{max width=\linewidth, max totalheight=0.42\textheight}{%
  \includegraphics{/Users/joschka/Desktop/Programming/eXercise/output/eXam/handouts/a_level_physics/assets/05-p104-positive-negative-work.pdf}}\\[0.3em]
\emph{Work done by the force (top) and work done on the force (bottom)}
\end{minipage}\par\medskip

For an object moving up a slope at angle $\alpha$ to the \textbf{horizontal} 水平, the work done against gravity in rising a height $h$ is $mgh$, while the work done by a horizontal push over the slope length $L$ uses $\cos\alpha$.

\subsection*{Conservation of energy}

Energy is \textbf{never made or destroyed} — it only changes from one form to another, or moves from one object to another. In a \textbf{closed system} 封闭系统, the total energy stays constant. This is \textbf{conservation of energy} 能量守恒.

When you write an energy equation, list every form the energy starts as and ends as. Common forms in this syllabus: kinetic, gravitational potential, \textbf{elastic potential energy} 弹性势能, \textbf{electrical} 电能, thermal, sound, \textbf{chemical energy} 化学能.

A ball rolling down a \textbf{frictionless} 无摩擦 \textbf{ramp} 斜坡 turns \textbf{gravitational potential energy} 重力势能 into \textbf{kinetic energy} 动能: $mgh = \tfrac{1}{2} m v^{2}$, so $v = \sqrt{2gh}$. With friction, some of this energy becomes \textbf{thermal energy} 热能 of the ramp and the air.

\subsection*{Efficiency}

The \textbf{efficiency} 效率 of a system is

\[
\text{efficiency} = \frac{\text{useful energy output}}{\text{total energy input}} \times 100\%.
\]

The same idea with power:

\[
\text{efficiency} = \frac{\text{useful power output}}{\text{total power input}} \times 100\%.
\]

Efficiency is always less than 100\% in a real system, because some input energy becomes "useless" forms — usually thermal energy.

For an electric motor lifting a load with efficiency $\eta$ at \textbf{voltage} 电压 $V$ and \textbf{current} 电流 $I$, the useful output power is $\eta V I$. From this you can find a force, a lifting speed, or a \textbf{tension} 张力.

\section*{Power}

\textbf{Power} 功率 is the rate of doing work, or the rate of transferring energy:

\[
P = \frac{W}{t} = \frac{\Delta E}{\Delta t}.
\]

Unit: $\text{W} = \text{J s}^{-1}$. Power is a scalar.

\subsection*{Power, force and velocity}

For an object moving at \textbf{velocity} 速度 $v$ with a force $F$ along the direction of motion, in a short time $\Delta t$ the displacement is $v\,\Delta t$ and the work done is $F v\,\Delta t$. Dividing by $\Delta t$:

\[
P = F v.
\]

This is one of the most useful results in mechanics.

\begin{itemize}
\item For a car at constant velocity $v$ on a flat road, the engine power must balance the total resistive force: $P = F_{\text{resist}} \cdot v$. If the \textbf{drag} 阻力 grows with $v^{2}$, doubling the speed roughly \textbf{quadruples} the power needed.

\item For lifting a \textbf{weight} 重力 $mg$ straight up at constant speed $v$, the useful output power is $P = mg \cdot v$.

\item For an aircraft hovering at a fixed height, the lift force equals the weight, and a large power is needed because air must be pushed downwards all the time.

\end{itemize}

\section*{Gravitational potential energy}

In a \textbf{uniform} gravitational field (close to a planet's surface), the change in gravitational potential energy of \textbf{mass} 质量 $m$ rising or falling through a height $\Delta h$ is

\[
\Delta E_{\text{P}} = m g \Delta h.
\]

\subsection*{Where it comes from}

The work done \textbf{against} gravity to raise a mass $m$ slowly (no change in kinetic energy) through height $\Delta h$ equals the gravitational potential energy gained:

\begin{itemize}
\item the gravitational force on the mass is $mg$ downwards,

\item the force needed to lift it slowly is $mg$ upwards,

\item the work done by this force is $W = F \cdot s = mg \cdot \Delta h$,

\item this work becomes $\Delta E_{\text{P}}$.

\end{itemize}

So $\Delta E_{\text{P}} = mg \Delta h$. To use it you need $m$ and $\Delta h$ (and $g$). You do not need speed or time.

\section*{Kinetic energy}

The kinetic energy of an object of mass $m$ moving at speed $v$ is

\[
E_{\text{K}} = \tfrac{1}{2} m v^{2}.
\]

\subsection*{Where it comes from}

Apply a resultant force $F$ to a mass $m$ that starts at rest. It speeds up evenly from $0$ to $v$ over a displacement $s$. From $v^{2} = u^{2} + 2as$ with $u = 0$,

\[
s = \frac{v^{2}}{2a}.
\]

The work done on the mass is

\[
W = F \cdot s = m a \cdot \frac{v^{2}}{2a} = \tfrac{1}{2} m v^{2}.
\]

All this work becomes kinetic energy, so $E_{\text{K}} = \tfrac{1}{2} m v^{2}$.

\subsection*{Kinetic energy and momentum}

Combining $p = mv$ and $E_{\text{K}} = \tfrac{1}{2} m v^{2}$:

\[
E_{\text{K}} = \frac{p^{2}}{2m}.
\]

This is handy when the \textbf{momentum} 动量 is given but not the velocity. For a momentum change from $p_{1}$ to $p_{2}$ at constant mass, the change in kinetic energy is $(p_{2}^{2} - p_{1}^{2}) / (2m)$.

\section*{Using energy methods}

A useful plan for problems that mix forces and energy:

\begin{enumerate}
\item \textbf{Find the start and end states.} Write the kinetic and potential energies in each.

\item \textbf{List any work done by outside forces} (friction, a push). Friction usually takes energy out; a push can add it.

\item \textbf{Conservation of energy}: $E_{\text{start}} + W_{\text{in}} = E_{\text{end}} + W_{\text{lost as heat etc.}}$.

\end{enumerate}

Examples:

\begin{itemize}
\item A box pushed at \textbf{constant velocity} up a ramp of length $L$ rising by $h$: $E_{\text{K}}$ does not change, so the work done by the push goes into $\Delta E_{\text{P}}$ plus the work done against friction.

\item A block sliding into a \textbf{spring} 弹簧 with kinetic energy $E_{\text{K}}$ on a frictionless surface: at greatest \textbf{compression} 压缩 $x$, all the kinetic energy has become elastic potential energy $\tfrac{1}{2} k x^{2}$ (where $k$ is the \textbf{spring constant} 劲度系数).

\item A ball dropped from height $h_{1}$ that bounces to height $h_{2}$: the ratio $h_{2}/h_{1}$ is the fraction of mechanical energy kept, $h_{2}/h_{1} = (v_{\text{up}}/v_{\text{down}})^{2}$.

\item A \textbf{projectile} 抛体 thrown to the same height at different angles: the final speed is the same (only the height matters); use components to get its direction.

\end{itemize}


\end{document}
